% !PRINTED 593
\weberpage{593}
%<<CONT
Componenten $x$ und $y$ absolut nicht gr\"osser als $1/2$ sind und dass also
\[
N(\xi-\mu) = x^2 - xy + y^2 \leqq \frac{3}{4}.
\]
Auf Grund dieser Eigenschaft l\"asst sich, ganz wie im vorigen Paragraphen, der gr\"osste gemeinschaftliche Theiler zweier Zahlen $\alpha$, $\beta$ bestimmen und damit beweisen, dass sich jede ganze Zahl in $R(\varrho)$ nur auf eine Weise in Primzahlen dieses K\"orpers zerlegen l\"asst, wenn associirte Zahlen nicht als wesentlich verschieden betrachtet werden.

Wie im vorigen Paragraphen schliesst man, dass man alle Primzahlen des K\"orpers $R(\varrho)$ erh\"alt, wenn man alle reellen Primzahlen zerlegt, und dass eine reelle Primzahl entweder in zwei conjugirte Primfactoren in $R(\varrho)$ zerf\"allt, $p=\pi\pi'$, oder dass sie auch im K\"orper $R(\varrho)$ eine Primzahl ist.

Die Primzahl 3 zerf\"allt in die beiden Factoren $(1-\varrho)$, $(1-\varrho^2)$; diese sind aber associirt und also ist $3=-\varrho^2(1-\varrho)^2$ mit dem Quadrat einer Primzahl in $R(\varrho)$ associirt. Man kann $\sqrt{-3}$ f\"ur diese Primzahl w\"ahlen. Alle Primzahlen von der Form $3f+1$ lassen sich nach \S.\,172 in zwei conjugirte Factoren in $R(\varrho)$ zerlegen, w\"ahrend die Primzahlen von der Form $3f+2$ niemals zerlegbar sind, weil ein Ausdruck von der Form $a^2-ab+b^2$ nicht $\equiv 2\ (\mathrm{mod}\ 3)$ sein kann. Also erh\"alt man alle Primzahlen in $R(\varrho)$, wenn man die reellen Primzahlen der Form $3f+2$ und die conjugirten Factoren der \"ubrigen reellen Primzahlen aufsucht.

Ist $a+b\varrho$ eine complexe ganze Zahl, deren Norm $a^2-ab+b^2$ nicht durch 3 theilbar ist, so muss von den beiden Zahlen $a$, $b$ entweder eine durch 3 theilbar sein, die andere nicht, oder es muss $a\equiv b$ und folglich $a^2-ab+b^2\equiv 1\ (\mathrm{mod}\ 3)$ sein. Es ist also unter den associirten Zahlen
\[
a+b\varrho,\quad b+a\varrho^2,\quad b-(a-b)\varrho,
\]
immer eine und nur eine, in der der Co\"efficient von $\varrho$ (oder von $\varrho^2$) durch 3 theilbar ist. Nehmen wir an, es sei $b=3B$ durch 3 theilbar, und setzen $N(a+b\varrho)=p$, so folgt, wenn noch $2a-b=A$ gesetzt wird,
\[
4p = A^2 + 27B^2,
\]
in Uebereinstimmung mit \S.\,172. Ist $p$ eine reelle Primzahl, und hat man $4p$ in dieser Form dargestellt, so erh\"alt man daraus die complexen Factoren von $p$
%CONT>>
